% ============================================
% Assistant Professor Cover Letter – St. Mary's College of Maryland (Economics / Finance)
% ============================================

\documentclass[11pt,letterpaper]{letter}

\usepackage{geometry}
\usepackage{microtype}
\usepackage[hidelinks]{hyperref}

\geometry{left=25mm,right=25mm,top=25mm,bottom=25mm}

\signature{Decory Edwards}
\address{
Department of Economics \\
Johns Hopkins University \\
Baltimore, MD 21218 \\
\texttt{dedwar65@jh.edu}
}

\begin{document}

\begin{letter}{Search Committee \\
Department of Economics \\
St. Mary’s College of Maryland \\
St. Mary’s City, MD 20686 \\ USA}

\opening{Dear Members of the Search Committee,}

I am writing to express my strong interest in the tenure-track Assistant Professor position in Economics at St. Mary’s College of Maryland. I am a Ph.D. candidate in Economics at Johns Hopkins University, advised by Professors Christopher D.~Carroll, Nicholas W.~Papageorge, and Stelios Fourakis, and I expect to receive my degree in May 2026. My research fields are macroeconomics, wealth and income dynamics, and computational economics.

My research agenda focuses on understanding the determinants of wealth inequality and the mechanisms through which heterogeneity in returns to wealth shapes aggregate outcomes. In my job-market paper, I estimate the degree of heterogeneity in individual portfolio returns using micro-data from the Survey of Consumer Finances and simulate a structural model in which households face idiosyncratic return risk. I show that allowing for persistent heterogeneity in returns substantially improves the model’s ability to match the empirical distribution of wealth, offering a tractable way to connect micro-level return dispersion to macroeconomic inequality.

In a second project, I use the Health and Retirement Study to examine the relationship between interpersonal trust and portfolio performance. Building on the literature that finds a hump-shaped relationship between trust and income, I show that a similar relationship holds for individual returns to wealth measured in the HRS data, highlighting a behavioral channel through which social attitudes shape saving and investment outcomes. This work also provides new estimates of the persistent component of returns to net worth, which motivate the calibration of heterogeneity in my structural model.

St. Mary’s College of Maryland is an especially attractive environment for this work. Its mission as Maryland’s public honors college, its emphasis on close faculty–student engagement, and its strong commitment to in-person teaching align closely with my approach to both research and instruction. I am particularly drawn to the department’s focus on finance-oriented courses such as Corporate Finance, Analyzing Financial Data, and Applied Portfolio Management, which naturally connect empirical analysis, theory, and real-world decision making. My research on portfolio returns and household financial behavior provides a strong foundation for engaging students with applied financial questions in a rigorous but accessible manner.

\textbf{Teaching.} I have experience teaching and supporting courses in \textit{Principles of Macroeconomics}, \textit{Intermediate Macroeconomic Theory}, and upper-level electives that emphasize data-driven economic analysis. My teaching philosophy emphasizes clarity, intuition, and active engagement, with a strong focus on connecting economic models to real financial decisions faced by households and firms. At St. Mary’s, I would be well prepared to teach Corporate Finance, Applied Portfolio Management, and courses focused on financial data analysis, while also contributing to the broader economics curriculum. I value inclusive, discussion-oriented classrooms and am committed to mentoring students from diverse backgrounds in both coursework and independent research.

Beyond academic fit, St. Mary’s location within Maryland and its proximity to Washington, D.C. make it an especially appealing place to build a long-term academic career. I am enthusiastic about contributing to the College’s mission through high-quality teaching, active scholarship, and sustained engagement with students and colleagues on campus.

Thank you very much for your time and consideration. I would welcome the opportunity to discuss my application further.

\closing{Sincerely,}

\end{letter}
\end{document}
